\documentclass[runningheads]{llncs}

\usepackage[T1]{fontenc}
\usepackage{graphicx}
\usepackage{url}
\usepackage{booktabs}

\begin{document}

\title{DevTrack AI: An Evidence-Based Traceability and Reporting System for Student Software Projects}

\titlerunning{DevTrack AI}

\author{Pham Duy Hung\inst{1} \and
Nguyen Thanh Dat\inst{1} \and
Nguyen Le Trung Tin\inst{1} \and
Tran Cong Tu\inst{1} \and
Nguyen Minh Hieu\inst{1}}

\authorrunning{N. M. Hieu et al.}

\institute{FPT University, Da Nang, Vietnam}

\maketitle

\begin{abstract}
Student software projects often combine requirement analysis, task implementation,
testing, source-code collaboration, and weekly reporting under strict time limits.
However, these activities are frequently managed through separated tools, making it
difficult to verify whether a requirement has been implemented, tested, supported by
evidence, and reflected in the final report. This paper presents DevTrack AI, a
project management support system designed for academic software teams. The system
organizes project artifacts through an end-to-end traceability flow from requirements
to use cases, tasks, test cases, defects, evidence, requirement traceability matrix
(RTM), and reports. DevTrack AI also introduces AI-assisted functions for artifact
suggestion and report drafting, while keeping official decisions under human review.
The proposed method combines rule-based traceability status calculation,
evidence-based task verification, GitHub metadata integration, and optional AI
summarization. The objective is to improve transparency, reduce missing artifacts,
and support fairer progress review in student project environments.

\keywords{Requirement traceability \and Agile project management \and Evidence-based verification \and Mining software repositories \and AI-assisted reporting}
\end{abstract}

% TODO: Add Jira link here if required by the instructor.
% Jira link: <PLACEHOLDER_FOR_JIRA_LINK>

\section{Introduction}

Software engineering courses require student teams to deliver a working product
together with requirements, tests, reports, and evidence of contribution. In
practice, these artifacts are often managed in separated tools. Requirements may be
written in documents, tasks may be tracked in a Kanban board, screenshots may be
stored in chat messages, test cases may be updated late, and code activity may stay
only in GitHub. This makes progress visible at the task level, but not always
verifiable at the requirement level. A leader or mentor may still be unable to
answer whether a requirement has been implemented, tested, supported by accepted
evidence, and included in a report.

Requirement traceability addresses this problem by connecting requirements with
downstream artifacts such as tasks, tests, source code, and change records
\cite{nlpTraceability,crossLevelTraceability}. However, maintaining traceability is
difficult in academic projects because teams work quickly, use lightweight tools,
and often prepare documentation near the end of a semester. At the same time, AI
tools can help draft tasks, tests, and reports, but AI-generated text is not reliable
unless it is grounded in structured project data. A fluent weekly report is not
enough if it cannot be traced back to requirements, evidence, and code activity.

DevTrack AI is proposed as an evidence-based traceability and reporting system for
student software teams. It is not designed as another task board. Instead, it
organizes project work around the artifact flow:

\begin{center}
\texttt{Requirement -> Use Case -> Task -> Test Case -> Bug/Defect -> Evidence -> RTM -> Report}
\end{center}

The goal is to make project progress explainable, not only visible. A task should
not be considered complete simply because its status is moved to \texttt{Done}; it
should be linked to a requirement, checked by tests or review evidence, and visible
in the Requirement Traceability Matrix (RTM). The RTM then becomes a live project
health view rather than a static table written only for final submission.

This paper contributes a lightweight system method for academic project tracking.
First, DevTrack AI uses rule-based RTM status calculation so that requirement health
is transparent. Second, it combines an Evidence Vault with a leader-controlled
review gate to reduce unsupported task completion. Third, it uses GitHub integration
and Code Insight to connect code evidence with project artifacts. Fourth, it uses AI
only as a support layer for suggestions and report drafts, while official decisions
remain under human review.

\section{Related Work}

\subsection{Agile Traceability and Requirement Traceability}

Requirement traceability connects requirements with related artifacts across the
software lifecycle. Prior work shows that natural language processing, information
retrieval, and learning-based techniques can help recover or maintain trace links
between textual artifacts \cite{nlpTraceability,crossLevelTraceability}. Other work
uses bag-of-words, TF-IDF, word embeddings, BERT-based models, and retrieval-augmented
generation to improve semantic matching and trace link recovery
\cite{crossLevelTraceability,traceBert,ragTraceability}.

These studies are useful for long-term automation, but they often focus on recovering
links after artifacts already exist. DevTrack AI takes a more operational direction
for student teams. It keeps traceability visible during project execution through a
rule-based RTM and human-confirmed artifact links. AI-based trace suggestions can be
added later, but official traceability evidence should remain explainable and
reviewable in the academic context.

\subsection{Requirement, Test, Defect, and Evidence Management}

Traceability is useful only when connected artifacts are meaningful and current. A
requirement should be related not only to implementation tasks, but also to acceptance
criteria, test cases, defects, and proof of completion. Requirement-driven tracking
therefore needs links between planning artifacts and verification artifacts so that
project status can be explained from evidence \cite{nlpTraceability,crossLevelTraceability}.

In student projects, this relationship is often weak. Teams may write test cases
late, record defects informally, or upload screenshots without review. DevTrack AI
addresses this gap through an Evidence Vault and a review-oriented task completion
flow. Evidence can include screenshots, API links, test results, GitHub issues,
commits, pull requests, and CI/check metadata. A task may request completion review,
but the leader still approves or rejects it based on the available evidence. This
keeps traceability useful for monitoring and collaboration, not only final
documentation \cite{ragTraceability}.

\subsection{Automated Reporting and Software Artifact Summarization}

Automated reporting is related to research on release note generation, pull request
description generation, and commit message generation. These studies commonly treat
software development data as source material for summarization. For example,
extractive summarization can select important commits or pull requests for release
notes \cite{releaseNotes}. Pull request description generation studies show how
commit messages and source-code context can be transformed into natural language
descriptions \cite{prDescriptions}. Commit message generation research similarly explores
how source-code changes can be summarized into readable descriptions
\cite{commitMessages}.

The lesson for DevTrack AI is that weekly reports should not be generated from an
empty prompt. The system first collects structured facts from tasks, sprint status,
RTM snapshots, defects, evidence review status, and GitHub activity. AI can then
rewrite those facts into a readable draft. This keeps the report grounded in project
data and allows the leader to review it before it is saved or submitted.

\subsection{Mining Software Repositories and AI for Software Engineering}

Mining Software Repositories (MSR) studies how data from version control systems,
issue trackers, pull requests, and other development platforms can be used to
understand software development activities \cite{teachingMsr}. This is relevant to
DevTrack AI because code contribution is an important part of student project
assessment, but commit count alone is not a fair indicator of contribution. DevTrack
AI therefore combines linked commits, pull requests, CI/check results, task
completion, accepted evidence, test activity, defect fixes, and review decisions.

Large language models (LLMs) have also been studied for software engineering tasks
such as code summarization, code understanding, issue analysis, and development
support \cite{llmSurvey}. In project-based
courses, LLM-based contribution summarization is also relevant because it can convert
repository activity into a more understandable description of individual work
\cite{llmContribution}. However, LLM output must be controlled because it can
produce plausible but unverifiable explanations. In DevTrack AI, AI support is
therefore used for explanation and drafting, not for automatic grading or task
approval. This human-in-the-loop principle is important for academic transparency.

\section{Method}

\subsection{Research and Design Approach}

This work follows a design-oriented approach. The target problem is identified from
common academic software project difficulties: unclear requirements, disconnected
tasks and tests, weak evidence, hard-to-review progress, and unclear individual
contribution. DevTrack AI is designed as a system artifact that combines
traceability management, evidence review, GitHub metadata integration, and
AI-assisted reporting.

The method does not claim that AI replaces project management decisions. Instead,
DevTrack AI separates deterministic project logic from AI support. Deterministic
logic is used for official status calculation, review gates, artifact relationships,
and evidence scoring. AI is used only for suggestions, structured review summaries,
and report drafts after project data has already been collected.

\subsection{System Artifact Model}

The core model follows the artifact chain introduced in the paper. A project contains
requirements. Requirements may have use cases and acceptance criteria. Tasks are
linked to requirements and assigned to members. Test cases verify requirements and
may create defects when execution fails. Evidence is uploaded or linked to tasks,
requirements, tests, defects, or code activity. The RTM aggregates these
relationships into a requirement-level health view, and reports are drafted from RTM
snapshots and project activity.

\begin{table}
\centering
\caption{Main DevTrack AI artifacts and their roles}
\begin{tabular}{ll}
\toprule
\textbf{Artifact} & \textbf{Role in the system} \\
\midrule
Requirement & Defines expected system behavior and acceptance criteria. \\
Use Case & Describes actor interaction and requirement scenarios. \\
Task & Represents implementation, testing, design, documentation, or review work. \\
Test Case & Verifies whether a requirement behaves as expected. \\
Bug/Defect & Records failed behavior found during test execution. \\
Evidence & Provides proof of implementation, testing, review, or demo result. \\
RTM & Summarizes traceability and requirement health. \\
Report & Presents weekly progress, risks, blockers, and contribution context. \\
\bottomrule
\end{tabular}
\end{table}

\subsection{Rule-Based RTM Status Calculation}

The RTM module is read-only over existing project data. It does not create or modify
requirements, tasks, tests, bugs, or evidence. Instead, it aggregates linked
artifacts and computes requirement health from task progress, test execution,
bug/defect state, accepted evidence, and risk reasons.

The status calculation is rule-based to keep the result explainable. A requirement is
\texttt{AT\_RISK} when it has blocking problems such as blocked tasks, failed or
blocked tests, open bugs, critical/high defects, or missing required evidence. It is
\texttt{DONE} when the requirement itself is done, linked tasks are done, linked
tests pass or are not blocking, and no blocking risk remains. It is
\texttt{NOT\_STARTED} when there is no task or test progress and no blocking risk.
Otherwise, it is \texttt{IN\_PROGRESS}. This makes the RTM a transparent review
surface: mentors can see not only the status, but also why the status was assigned.

\subsection{Evidence-Based Task Review}

DevTrack AI uses a review gate before a task is finally accepted as done. A member
requests review after completing work. The system builds an evidence summary from
available project data, including requirement linkage, assignee information,
checklist status, subtasks, blocked state, linked GitHub issue, uploaded evidence,
and linked code evidence. The project leader then approves or rejects the task. This
process reduces false completion, where a task is marked done without meaningful
proof.

The Code Insight module extends this review gate for code-related tasks. It uses the
shared GitHub Integration Core to handle repository configuration, webhook
verification, and event dispatching. Push, pull request, workflow, and check-run
events can be stored as normalized GitHub evidence and linked to tasks through
explicit task identifiers, GitHub issue references, or SHA chains. Code Insight then
uses local task signals and linked GitHub evidence to produce an explainable score,
show evidence details to the leader, fetch changed-file metadata when needed, store
local structured AI review summaries, capture review snapshots, and expose dashboard
metrics. The leader remains the final authority:

\begin{center}
\texttt{GitHub evidence + deterministic rules + optional AI review + leader decision = trusted task completion}
\end{center}

\subsection{GitHub Integration and Code Insight}

GitHub integration is treated as shared infrastructure rather than a repository
clone. DevTrack AI stores only the metadata needed for traceability, review,
reporting, and audit. This includes repository configuration, webhook events,
commits, pull requests, CI/check results, changed-file summaries, evidence links,
review decisions, and Code Insight snapshots.

This design supports evidence-based project review without duplicating GitHub. A
task can be checked for linked commits, pull requests, failed CI, weak evidence, or
missing code activity. Because the system keeps the evidence and the leader decision
together, source-code activity becomes part of the project traceability model rather
than a separate repository history.

\subsection{AI-Assisted Reporting}

The weekly report method uses structured project facts as input. The system collects
task status, sprint scope, requirement status, RTM snapshot, test execution results,
defects, evidence review status, and GitHub activity. These facts are then grouped
into progress, completed work, technical changes, risks, blockers, and contribution
context. AI may be used to rewrite this structured summary into a readable report
draft.

The leader reviews the report before it is stored or submitted. This step is
important because AI-generated language may hide missing data if the input is
incomplete. In DevTrack AI, the AI report is therefore a draft, not the official
source of truth. The official source remains the structured project data, accepted
evidence, RTM state, and review decisions.

\subsection{Evaluation Direction}

The planned evaluation focuses on traceability completeness and review usefulness.
Student project scenarios can be used to check whether DevTrack AI identifies
requirements without tasks, tasks without evidence, requirements without test cases,
failed tests without defects, unlinked code evidence, and weak completion claims. For
AI-assisted reporting, generated drafts should be compared against structured project
facts to detect unsupported claims. Feedback from leaders, members, and mentors can
also be collected to assess whether the RTM and evidence workflow make project
review clearer.

\section{References}

\begin{thebibliography}{99}

\bibitem{nlpTraceability}
Guo, J.L.C., Steghöfer, J.P., Vogelsang, A., Cleland-Huang, J.: Natural Language 
Processing for Requirements Traceability. arXiv preprint arXiv:2405.10845 (2024).

\bibitem{ragTraceability}
Hey, T., Fuchß, D., Keim, J., Koziolek, A.: Requirements Traceability Link 
Recovery via Retrieval-Augmented Generation. In: International Working 
Conference on Requirements Engineering: Foundation for Software Quality (REFSQ). 
Lecture Notes in Computer Science. Springer, Cham (2025).

\bibitem{crossLevelTraceability}
Mohammad, B., Sonbol, R., Rebdawi, G.: Cross-level Requirement Traceability: 
A Novel Approach Integrating Bag-of-Words and Word Embedding for Enhanced 
Similarity Functionality. arXiv preprint arXiv:2406.14310 (2024).

\bibitem{traceBert}
Lin, J., Liu, Y., Zeng, Q., Jiang, M., Cleland-Huang, J.: Traceability 
Transformed: Generating More Accurate Links with Pre-Trained BERT Models. 
arXiv preprint arXiv:2102.04411 (2021).

\bibitem{releaseNotes}
Nath, S.S., Roy, B.: Towards Automatically Generating Release Notes using 
Extractive Summarization Technique. arXiv preprint arXiv:2204.05345 (2022).

\bibitem{prDescriptions}
Liu, Z., Xia, X., Treude, C., Lo, D., Li, S.: Automatic Generation of Pull 
Request Descriptions. arXiv preprint arXiv:1909.06987 (2019).

\bibitem{commitMessages}
Cortés-Coy, L.F., Linares-Vásquez, M., Aponte, J., Poshyvanyk, D.: On
Automatically Generating Commit Messages via Summarization of Source Code Changes.
In: Proceedings of the 14th IEEE International Working Conference on Source Code
Analysis and Manipulation (SCAM), pp. 275--284 (2014).

\bibitem{teachingMsr}
Codabux, Z., Fard, F., Verdecchia, R., Palomba, F., Di Nucci, D., Recupito, G.: 
Teaching Mining Software Repositories. arXiv preprint arXiv:2501.01903 (2025).

\bibitem{llmSurvey}
Zhang, Q., Fang, C., Xie, Y., Zhang, Y., Yang, Y., Sun, W., Yu, S., Chen, Z.: 
A Survey on Large Language Models for Software Engineering. arXiv preprint 
arXiv:2312.15223 (2024).

\bibitem{llmContribution}
Ferrao, R.C., de Miranda, F., Soler, D.P.: LLM Contribution Summarization 
in Software Projects. arXiv preprint arXiv:2505.17710 (2025).

\end{thebibliography}

\end{document}
